\chapter{Conclusions}
This thesis has presented the design, implementation, and experimental evaluation of a cascaded deep learning architecture for Vietnamese optical character recognition, combining a fine-tuned RT-DETR V2 oriented text line detector with a fine-tuned PARSeq permutation-based sequence recognizer. The system was developed in response to two interconnected gaps identified in Chapter \ref{chapter:related_works}: the absence of a Vietnamese OCR pipeline capable of handling arbitrarily oriented text line detection with tight bounding box localization, and the limited performance of existing recognition systems on the morphologically complex Vietnamese character set. The following sections summarize the principal contributions of the work, reflect on how the proposed system addresses the challenges outlined in the introduction, and draw conclusions from the experimental findings.

\section{Summary of Contributions}

The thesis makes four principal technical contributions, each addressing a distinct bottleneck in the Vietnamese OCR pipeline.

The first contribution is the extension of RT-DETR V2 with an auxiliary oriented bounding box detection capability. The standard RT-DETR decoder, which produces axis-aligned box predictions, was augmented with a per-layer angle regression head implemented as a three-layer MLP with sine-cosine output parameterization. This extension enables the detector to predict text line orientations across the full $[0^\circ, 360^\circ)$ range without the angular discontinuity that afflicts scalar angle regression, while the GWD loss \cite{Yang2022} provides geometrically unified supervision over all five oriented box parameters jointly. The combination of sine-cosine encoding and GWD loss constitutes the primary methodological contribution to the detection stage and is directly motivated by the observation, formalized in Section \ref{sec:error_propagation_analysis_across_stages}, that angular localization errors in the detection stage disproportionately degrade Vietnamese diacritical recognition by displacing tone marks from the receptive fields of the ViT patch embeddings in the second stage.

The second contribution is the development of a large-scale synthetic data generation pipeline specifically designed for Vietnamese oriented text detection training. The pipeline produces approximately 130,000 annotated detection images spanning seven document layout archetypes, including single-column administrative documents, receipts, and fully rotated scatter layouts, with exact oriented bounding box annotations derived analytically from rendering parameters. A critical engineering contribution within this pipeline is the targeted patch that corrects the angle distribution collapse caused by the use of \texttt{cv2.minAreaRect} in the base generator, replacing it with a direct PIL-based rendering and analytical annotation approach that produces a flat angle distribution across the full rotation range. This correction was essential for enabling the angle regression head to generalize beyond the degenerate near-horizontal distribution that the uncorrected generator produced.

The third contribution is the adaptation of PARSeq to the Vietnamese recognition domain through vocabulary extension, class-weighted training, and phonological validity constraints. The output vocabulary was expanded from 94 ASCII characters to approximately 310 tokens covering the full set of Vietnamese composite characters, with warm-start weight initialization from the pretrained checkpoint for shared characters and random initialization for new Vietnamese glyphs. Class-weighted cross-entropy loss was applied to counteract the severe frequency imbalance between common unaccented characters and rare tone-diacritic combinations. At inference time, a phonological validity mask suppresses outputs corresponding to impossible Vietnamese character combinations, reducing the rate of linguistically invalid predictions under degraded image conditions. These adaptations, combined with fine-tuning on a 2,000,000-image synthetic recognition corpus generated by a three-branch parallel generation pipeline, produced a recognizer that outperforms all evaluated baselines on the Vietnamese test set.

The fourth contribution is the system integration and modular pipeline architecture that connects the two model components through a well-defined inter-stage interface. The interface performs perspective-corrected cropping using the predicted oriented bounding box parameters, converting arbitrarily oriented text line regions to upright rectangular crops before passing them to the recognizer. This cropping strategy, detailed in Section \ref{sec:inter_stage_data_flow_and_cropping_strategy_3_1_2}, directly operationalizes the theoretical motivation for oriented detection: tighter angular localization reduces background inclusion in inter-stage crops and thereby reduces the recognition error rate on diacritically complex Vietnamese words.

\section{Addressing the Challenges Outlined in the Introduction}

Chapter 1 identified three principal challenges for Vietnamese OCR: the complexity of the Vietnamese diacritical system, the diversity of Vietnamese document formats and text orientations, and the scarcity of annotated Vietnamese OCR data for training and evaluation.

The diacritical complexity challenge is addressed at both stages of the pipeline. In the detection stage, the oriented bounding box representation and GWD loss ensure that text line crops are delivered to the recognizer with minimal residual rotation, which is critical because even small angular errors displace diacritical marks out of the ViT patch embeddings that encode their corresponding base characters. In the recognition stage, the extended Vietnamese vocabulary, class-weighted loss, and phonological validity constraints collectively ensure that the model can represent, learn, and constrain its outputs to the full Vietnamese composite character inventory. The experimental results confirm that this combination achieves a CER of 10.18\% and WER of 34.71\%, outperforming all evaluated baselines including VietOCR, which was specifically designed for Vietnamese text.

The document format diversity challenge is addressed through the multi-layout synthetic generation pipeline, which produces training images spanning seven document archetypes covering the range from structured administrative correspondence to arbitrarily oriented scatter text. The rotated-scatter layout variant specifically ensures that the angle regression head receives training signal across the full $[0^\circ, 360^\circ)$ range, enabling the detector to handle document images with any text orientation without requiring layout-specific model variants.

The data scarcity challenge is addressed through fully synthetic data generation, which eliminates the need for manual annotation effort while producing training sets of sufficient scale for both detection (130,000 images) and recognition (2,000,000 crops). The augmentation pipelines applied to both corpora introduce controlled photometric and geometric variation that partially bridges the gap between synthetic and real document conditions, though the experimental results indicate that a residual domain gap remains and constitutes a primary direction for future improvement.

\section{Experimental Conclusions}

The experimental evaluation in Chapter \ref{chapter:experiments_4} confirms the effectiveness of the proposed system along both dimensions of the OCR pipeline. On the detection task, the fine-tuned RT-DETR model achieves F1@0.50 of 0.501 and mAP@0.5:0.95 of 0.147, compared to 0.134 and 0.008 for the EAST baseline, representing substantial improvements across all reported metrics. The disproportionate advantage at stricter IoU thresholds, where the RT-DETR model achieves F1@0.75 of 0.188 versus 0.021 for EAST, confirms that the oriented detection head and GWD loss contribute directly to localization tightness beyond what axis-aligned detection can achieve. On the recognition task, the fine-tuned PARSeq model achieves the best CER and WER among all five evaluated systems, with improvements over the most competitive baseline VietOCR of 1.16 percentage points in CER and 2.73 percentage points in WER.

The remaining performance gap, reflected in the WER of 34.71\%, indicates that Vietnamese OCR remains an open and challenging problem. The primary bottlenecks identified through error analysis are the visual ambiguity between diacritically similar character pairs under degraded image conditions, and the domain gap between the synthetic training distribution and real-world scanned documents. Both of these limitations are addressable through future work involving real data collection and annotation, higher-resolution recognition processing, and joint optimization of the detection and recognition stages.

\section{Practical Value for Vietnamese Document Digitization}

The practical value of the proposed system lies in its modular, adaptable architecture and its demonstrated superiority over general-purpose OCR tools on Vietnamese text. Vietnamese administrative agencies, archives, and enterprises generate large volumes of printed documents that require digitization for indexing, information extraction, and long-term preservation. Existing general-purpose OCR systems, as demonstrated by the baseline results in Chapter \ref{chapter:experiments_4}, are poorly suited to Vietnamese text due to their limited character set coverage and lack of Vietnamese-specific optimization. The proposed system addresses this gap by providing a complete, end-to-end pipeline that can be deployed through a single command-line interface, produces structured Markdown transcripts alongside visualization outputs, and can be adapted to new document types or updated language models by replacing individual components without retraining the entire system.

The cascaded modular design is particularly well suited to the administrative document digitization context, where document format specifications change over time and different deployment scenarios may require different trade-offs between detection accuracy and inference speed. The independent encapsulation of the detector in the \texttt{LineDetector} class and the recognizer in the \texttt{TextRecognizer} class allows either component to be upgraded or fine-tuned independently as improved model checkpoints become available, without disrupting the pipeline interface.

In summary, this thesis demonstrates that a carefully designed combination of domain-specific synthetic data generation, oriented bounding box detection with geometrically principled loss functions, and Vietnamese-adapted permutation-based recognition produces a Vietnamese OCR system that substantially outperforms the available alternatives and addresses the core technical challenges that have limited the effectiveness of general-purpose OCR on Vietnamese text. The system represents a practical step toward scalable, automated Vietnamese document digitization and establishes a foundation for the future improvements identified in Chapter \ref{chapter:experiments_4}.
